% Removed: Word title page, placeholder Abstract template text, and Word-generated indices (TOC/figures/tables/abbrev lists).
% Kept below: the actual theory text as editable LaTeX.

\section{\texorpdfstring{\textbf{Background and
Rationale}}{Background and Rationale}}\label{background-and-rationale}

Uzbekistan's health system, like those of many post-Soviet states, has
undergone a complex and protracted reform process since independence in
1991. The legacy of the Semashko model - characterized by centralized
planning, public financing and provision, and a strong hospital sector -
has shaped both the strengths and persistent challenges of the system.
Over the past three decades, the country has sought to replace
input-driven norms with outcomes-focused governance, strengthen primary
health care (PHC), and reduce out-of-pocket (OOP) payments that threaten
financial protection and equity. (Ketevan Glonti, 2013)

Despite important gains - particularly in routine immunization, maternal
and child health, and the introduction of strategic purchasing pilots -
Uzbekistan continues to face high OOP payments, gaps in financial
protection, underdeveloped quality assurance, and pronounced
rural--urban disparities. The burden of non-communicable diseases (NCDs)
has risen sharply, with ischemic heart disease, stroke, and cancers now
the leading causes of mortality. These epidemiological pressures,
combined with infrastructure and workforce constraints, have exposed the
limitations of traditional, hospital-centric models and underscored the
need for innovative, data-driven, and patient-centered approaches.
(World Health Organization, 2022)

In this context, digital health---including telemedicine platforms,
electronic health records, and e-referral systems---has emerged as a
promising instrument for health system improvement. The government's
``Concept on Health Development 2019--2025'' and subsequent strategies
have prioritized digitalization as a lever for expanding access,
improving quality, and enhancing efficiency. Early pilots in Syrdarya
region, including telemedicine clinics, e-referral systems, and digital
queue management, represent a notable inflection point in the country's
reform trajectory. (World Health Organization, 2022)

However, the digital health agenda is nascent and faces significant
barriers: legacy paper-based systems, fragmented data flows, uneven
facility readiness, limited digital literacy, and the absence of robust
governance and interoperability standards. Against this backdrop, this
thesis offers an integrated theoretical examination of health system
transformation in Uzbekistan, anchored in the specificities of the Uzbek
case and informed by comparative regional and international experience.

\section{\texorpdfstring{\textbf{Problem Statement and Research
Questions}}{Problem Statement and Research Questions}}\label{problem-statement-and-research-questions}

\textbf{Problem Statement:}

Uzbekistan's health system faces persistent challenges in healthcare
access, service efficiency, and equity - particularly in rural areas and
among vulnerable populations. Despite ongoing reforms and digital health
pilots, gaps remain in infrastructure, coordination, and responsiveness.
Telemedicine platforms offer a promising solution, but their design must
be grounded in the lived experiences and needs of both patients and
healthcare providers. This thesis investigates how a telemedicine
platform can be developed and tested to address these challenges, based
on direct input from users and providers within the Uzbek health system.

\textbf{Research Questions:}

\begin{itemize}
\item
  What are the main barriers to healthcare access as perceived by
  patients in Uzbekistan, and how can digital tools address these
  barriers?
\item
  What are the key service delivery inefficiencies experienced by
  doctors, and what features would improve their workflow and patient
  management through a telemedicine platform?
\item
  What functional and technical requirements should a telemedicine
  platform meet to be usable, acceptable, and effective for both
  patients and providers in Uzbekistan?
\item
  How do users (patients and doctors) respond to a prototype
  telemedicine platform, and what does their feedback reveal about its
  potential to improve access and efficiency?
\end{itemize}

\section{\texorpdfstring{\textbf{1.3 Objectives of the
Thesis}}{1.3 Objectives of the Thesis}}\label{objectives-of-the-thesis}

\begin{itemize}
\item
  To identify key barriers to healthcare access and service delivery
  efficiency in Uzbekistan through empirical surveys of patients and
  healthcare providers.
\item
  To translate user-reported challenges into functional and technical
  requirements for a telemedicine platform tailored to Uzbekistan's
  health system context.
\item
  To design and develop a prototype digital portal for patients and
  doctors that addresses the identified needs and supports improved
  access and workflow efficiency.
\item
  To evaluate the usability, acceptability, and perceived effectiveness
  of the telemedicine platform through pilot testing with real users.
\item
  To contribute to the theoretical and practical understanding of
  user-centered digital health design in post-Soviet health systems
  undergoing reform.
\end{itemize}

Healthcare System in Uzbekistan

\section{2.1 Historical Background and Current
Organization}\label{historical-background-and-current-organization}

Uzbekistan's health system is a product of its Soviet legacy, shaped by
the Semashko model---a highly centralized, state-funded, and
state-delivered system emphasizing hospital-based care, vertical disease
programs, and quantitative targets for inputs such as beds and staff
(Bernd Rechel M. A., 2014). During the Soviet era, health care was
nominally free at the point of use, with a focus on communicable disease
control, maternal and child health, and the expansion of physical
infrastructure. However, the system was characterized by inefficiencies,
underfunding, and a neglect of primary care and quality assurance (Bernd
Rechel E. R., 2014).

Following independence in 1991, Uzbekistan faced the dual challenge of
economic transition and health system reform. The early 1990s saw a
sharp contraction in public spending, a decline in health indicators,
and the emergence of informal payments and private provision (Bernd
Rechel M. A., 2014). The government responded with a series of reforms
aimed at rationalizing hospital capacity, strengthening PHC, and
introducing new financing mechanisms.

\textbf{Key Reform Milestones:}

\begin{itemize}
\item
  \textbf{1996 Law on Health Protection:} Established the legal basis
  for a state-guaranteed benefits package, including primary care,
  emergency care, and care for ``socially significant and hazardous''
  conditions (Bernd Rechel M. A., 2014).
\item
  \textbf{1998 Presidential Decree:} Launched a state program for health
  system reform, prioritizing PHC, emergency care, and the development
  of the private sector (Bernd Rechel M. A., 2014).
\item
  \textbf{2000s--2010s:} Implementation of World Bank and ADB-supported
  projects (Health I, II, III; Woman and Child Health Development)
  focused on PHC infrastructure, maternal and child health, and
  financing reforms (Bernd Rechel M. A., 2014). (The World Bank, 2025)
\item
  \textbf{2018--2021:} Adoption of the ``Concept on Health Development
  2019--2025'' (Presidential Decree 5590), establishment of the SHIF,
  and piloting of strategic purchasing and digital health initiatives in
  Syrdarya (The World Bank, 2025).
\end{itemize}

\textbf{Current Organization:}

The health system is organized into three administrative tiers:

\begin{itemize}
\item
  \textbf{National (Republican) Level:} Ministry of Health, national
  research institutes, tertiary hospitals, and regulatory agencies.
\item
  \textbf{Regional (Viloyat) Level:} Regional health authorities,
  multi-specialty hospitals, and specialized centers.
\item
  \textbf{District/City Level:} District medical unions, central
  hospitals, family polyclinics, and rural physician points (Bernd
  Rechel M. A., 2014).
\end{itemize}

The MoH remains the principal steward, with direct management of
national-level institutions and regulatory oversight of regional and
district facilities. The private sector, while growing, is still limited
in scope and subject to strict licensing and regulatory controls (The
World Bank, 2025).

\textbf{Service Delivery Model:}

\begin{itemize}
\item
  \textbf{Primary Care:} Delivered through family doctor points,
  polyclinics, and rural physician posts. PHC reforms have aimed to
  shift from specialist-dominated polyclinics to generalist-led family
  medicine, but progress has been uneven (Bernd Rechel A. S., 2023).
\item
  \textbf{Secondary and Tertiary Care:} Provided by district, regional,
  and national hospitals, with a legacy of overcapacity and
  fragmentation. Recent reforms aim to consolidate specialized centers
  into multi-specialty general hospitals (The World Bank, 2025)
\item
  \textbf{Emergency Care:} A well-developed network of emergency
  departments often better equipped than other public facilities, but
  subject to overload due to gaps in PHC and benefits coverage (World
  Health Organization, 2022).
\end{itemize}

\textbf{Governance and Regulation:}

Regulation is highly centralized, with the MoH and Cabinet of Ministers
responsible for policy, financing, and oversight. Local authorities have
limited autonomy, and dual accountability structures persist (Bernd
Rechel M. A., 2014). The regulatory framework for private providers and
pharmaceuticals is evolving, with recent efforts to strengthen
licensing, quality assurance, and market surveillance (The World Bank,
2025).

\section{\texorpdfstring{\textbf{2.2 Key Challenges in Access, Quality,
and
Financing}}{2.2 Key Challenges in Access, Quality, and Financing}}\label{key-challenges-in-access-quality-and-financing}

\textbf{Access:}

Despite improvements in health indicators, access to quality care
remains uneven. Rural-urban disparities are pronounced, with rural areas
facing shortages of qualified staff, inadequate infrastructure, and
limited access to diagnostics and medicines (Bernd Rechel A. S., 2023).
For example, only 57\% of PHC facilities reported basic water services
and 26\% basic sanitation in 2020, reflecting broader readiness
constraints (World Health Organization, 2022).

Geographical access is further constrained by the distribution of
facilities and workforce. While the number of physicians per 100,000
population has increased to 276 by 2019, and nurse density to
\textasciitilde1,100 per 100,000, these resources are concentrated in
urban centers. Outward migration of health professionals, especially to
Russia and Kazakhstan, exacerbates shortages in underserved areas.
(Bernd Rechel E. R., 2014)

\textbf{Quality:}

Quality of care is undermined by several factors:

\begin{itemize}
\item
  \textbf{Legacy of Input-Based Planning:} The Soviet system prioritized
  quantitative targets (beds, staff) over outcomes, with limited
  attention to quality assurance, patient safety, or evidence-based
  practice (Ketevan Glonti, 2013).
\item
  \textbf{Outdated Clinical Protocols:} Many facilities rely on outdated
  or poorly adapted clinical guidelines, with limited adoption of
  international best practices (Bernd Rechel M. A., 2014).
\item
  \textbf{Weak Monitoring and Evaluation:} Data systems are fragmented,
  paper-based, and focused on structural indicators rather than process
  or outcome measures. Quality improvement initiatives are limited in
  scope and sustainability (The World Bank, 2025)
\item
  \textbf{Patient Experience:} User experience is often neglected, with
  long waiting times, limited appointment systems, and a lack of
  patient-centered care models (Bernd Rechel M. A., 2014).
\end{itemize}

\textbf{Financing:}

The health financing system is characterized by:

\begin{itemize}
\item
  \textbf{Low Public Spending:} Public spending on health was
  \textasciitilde2.3\% of GDP in 2019, above the Central Asia average
  but well below the WHO European Region average (\textasciitilde5\%)
  (World Health Organization, 2022).
\item
  \textbf{High Out-of-Pocket Payments:} OOP payments accounted for
  57.7\% of current health expenditure in 2019, exposing households to
  financial hardship and catastrophic health spending. In 2020, 18\% of
  households reported cost-related nonuse of care (World Health
  Organization, 2022)
\item
  \textbf{Fragmented Pooling and Purchasing:} Budgeting is historically
  incremental, with limited strategic purchasing or risk pooling. The
  SHIF pilot in Syrdarya aims to introduce capitation and DRG-based
  payments, but implementation is nascent and faces capacity constraints
  (The World Bank, 2025).
\item
  \textbf{Limited Benefits Coverage:} The basic benefits package covers
  primary and emergency care, but excludes much of secondary/tertiary
  care and outpatient medicines, amplifying OOP and incentivizing
  emergency use (The World Bank, 2025).
\end{itemize}

\textbf{Equity and Financial Protection} (Bernd Rechel M. A.,
2014)\textbf{:}

\begin{itemize}
\item
  \textbf{Vulnerable Groups:} While the benefits package includes
  provisions for vulnerable groups (e.g., children, pregnant women,
  people with disabilities), eligibility criteria are not directly
  linked to income, and coverage gaps persist.
\item
  \textbf{Informal Payments:} Informal payments remain widespread,
  particularly in secondary and tertiary care, undermining equity and
  trust in the system.
\end{itemize}

\section{\texorpdfstring{\textbf{2.3 Health Outcomes and Risk
Profile}}{2.3 Health Outcomes and Risk Profile}}\label{health-outcomes-and-risk-profile}

\textbf{Demographic and Epidemiological Trends} (WHO Regional Office for
Europe, 2023)\textbf{:}

\begin{itemize}
\item
  \textbf{Population:} 33.3 million in 2020, with a young age structure
  but a growing burden of NCDs.
\item
  \textbf{Life Expectancy:} Increased from 67.2 years in 2000 to 73.9
  years in 2016, with a gender gap of 4.7 years.
\item
  \textbf{Maternal and Infant Mortality:} Maternal mortality declined
  from 41 per 100,000 live births in 2000 to 29 in 2017; infant
  mortality from 51.8 per 1,000 in 2000 to 15.6 in 2019.
\item
  \textbf{NCDs:} Account for 84\% of all deaths in 2018, with ischemic
  heart disease and stroke as leading causes.
\item
  \textbf{Risk Factors:} High systolic blood pressure (30.7\%), dietary
  risks (28.2\%), and high fasting plasma glucose (20.2\%) are dominant
  contributors to mortality.
\end{itemize}

\textbf{Service Coverage and UHC Index} (World Health Organization,
2022)\textbf{:}

\begin{itemize}
\item
  \textbf{UHC Service Coverage Index:} Improved from 56 (2000) to 71
  (2019), below OECD/EU comparators but signaling progress.
\item
  \textbf{Immunization:} High coverage rates for childhood vaccines, but
  challenges remain in TB, HIV/AIDS, and MDR-TB.
\end{itemize}

\textbf{Social Determinants} (The World Bank Group, 2021)\textbf{:}

\begin{itemize}
\item
  \textbf{Poverty:} Poverty headcount ratio declined from 14.1\% in 2013
  to 11.4\% in 2018, but disparities persist.
\item
  \textbf{Water and Sanitation:} Infrastructure needs extensive
  rehabilitation, with poor access in many areas.
\end{itemize}

Digitalization \& Internet Infrastructure

\section{\texorpdfstring{\textbf{3.1 Introduction: Digital Health as a
Systemic
Enabler}}{3.1 Introduction: Digital Health as a Systemic Enabler}}\label{introduction-digital-health-as-a-systemic-enabler}

Digital transformation is now recognized as a foundational enabler of
health system reform, not merely a technical add-on. In Uzbekistan, the
digital health agenda is explicitly linked to the broader goals of
universal health coverage (UHC), equity, efficiency, and quality
improvement. The government's ``Digital Uzbekistan 2030'' strategy and
the National Health System Strategy 2030 both position digitalization as
a cross-cutting priority, aiming to create robust governance structures,
integrated platforms, and the necessary broadband and ICT infrastructure
to support health system modernization. (The World Bank, 2025)

However, the digital health landscape in Uzbekistan is shaped by the
legacy of Soviet-era information systems, the rapid but uneven expansion
of internet and mobile technologies, and persistent gaps in digital
literacy and organizational readiness. This section provides a
comprehensive analysis of the current state, challenges, and
opportunities for digitalization and internet infrastructure in
Uzbekistan's health sector.

\section{\texorpdfstring{\textbf{3.2 Post-2016 Reforms: Infrastructure,
Connectivity, and Policy} (The World Bank,
2025)}{3.2 Post-2016 Reforms: Infrastructure, Connectivity, and Policy (The World Bank, 2025)}}\label{post-2016-reforms-infrastructure-connectivity-and-policy-the-world-bank-2025}

Since 2016, Uzbekistan has embarked on an ambitious digital
transformation agenda, with health as a priority sector. The ``Digital
Uzbekistan 2030'' strategy sets out to strengthen digital governance,
create integrated information systems, and expand broadband connectivity
across the country. The National Health System Strategy 2030 further
articulates the role of digital health as a driver of system-wide
reform, emphasizing the need for interoperable platforms, data-driven
decision-making, and patient empowerment.

By 2023, up to 81\% of public medical organizations' legal entities had
been provided with internet connectivity. However, the quality and
reliability of internet access vary significantly, especially in rural
areas. The phased installation of local area networks (LANs) in
healthcare facilities is ongoing. Despite these advances, more than 50\%
of public medical organizations still lack LANs beyond administrative
offices, and the estimated need for computers ranges from 30,000 to
65,000 units.

The lack of adequate ICT infrastructure is a critical bottleneck for
digital health. Many facilities are underequipped, with insufficient
computers, outdated hardware, and no dedicated ICT maintenance staff.
The absence of robust data centers further limits the scalability of
national digital health systems.

\section{\texorpdfstring{\textbf{3.3 National and Local Health
Information
Systems}}{3.3 National and Local Health Information Systems}}\label{national-and-local-health-information-systems}

Uzbekistan's health sector is characterized by a proliferation of more
than 35 unconnected nationwide IT e-health systems. These systems
operate in silos, with no integration or common master data usage among
Ministry of Health (MoH) information systems. There is no national
patient registry, no unified registry of healthcare providers, and no
standardized nomenclature or classifiers. As a result, integration of IT
systems across facilities is not feasible, and data exchange is severely
constrained. (WHO Regional Office for Europe, 2023)

At the facility level, only a limited number of organizations use
Patient Administration Systems (PAS), Electronic Medical Records (EMR),
or Hospital Information Systems (HIS). Under MoH initiatives, two
PAS/EMR systems have been piloted in Syrdarya region, and about five
local PAS/HIS vendors provide services to public and private clinics and
hospitals. However, these systems are not interoperable, and most
medical records remain paper-based, scattered across facilities, and
inaccessible for longitudinal patient care. (The World Bank, 2025)

\section{\texorpdfstring{\textbf{3.4 Digital Litracy and Adoption Rates}
(WHO Regional Office for Europe,
2023)}{3.4 Digital Litracy and Adoption Rates (WHO Regional Office for Europe, 2023)}}\label{digital-litracy-and-adoption-rates-who-regional-office-for-europe-2023}

Digital literacy among health professionals is highly variable. While
younger staff and those in urban centers are increasingly proficient
with computers and digital tools, many older providers lack basic ICT
skills. Continuous IT training is not yet institutionalized, and there
is no national strategy for digital capacity building in the health
workforce. This gap is compounded by the lack of standardized digital
workflows and the persistence of paper-based documentation, which
consumes significant provider time and undermines efficiency.

Patient digital literacy is also uneven, with urban populations more
likely to use digital health services than rural residents. The absence
of a national patient portal or unified engagement platform limits
opportunities for patients to access their health records, book
appointments, or participate in their care. Where electronic booking
systems exist, their use is limited and not widely spread. As a result,
patients are not ``in the driver's seat'' of their health journey, and
digital health remains largely provider-centric.

The success of digital health initiatives depends not only on technical
infrastructure but also on organizational readiness and change
management. In Uzbekistan, digitalization planning and management are
currently centralized within the MoH and its IT-focused company,
UZINFOCOM. There is no independent body tasked with overseeing digital
health development, and the process is often non-transparent to external
stakeholders, including hospitals, clinics, professional associations,
and patient groups. This lack of inclusive governance hinders buy-in and
slows the pace of digital transformation. (The World Bank, 2025)

Literaturverzeichnis

Bernd Rechel, E. R. (2014). Trends in health systems in the former
Soviet countries. \emph{European Journal of Publich Health}.

Ketevan Glonti, B. R. (2013). Health Targets in the Former Soviet
Countries: Responding to the NCD Challenge? \emph{Public Health Reviews,
Vol. 35, No 1}.

The World Bank. (2016). \emph{SYSTEMATIC COUNTRY DIAGNOSTIC for
UZBEKISTAN.} Washington D.C: The World Bank.

World Health Organization. (2022). \emph{Health Systems in Action
Uzbekistan.} World Health Organization.

Anhang

Füge hier alle Anlagen hinzu, auf die du in deiner Arbeit verweist.

Du kannst die lesende Person im Fließtext auf Dokumente im Anhang
aufmerksam machen, indem du nach einer Aussage in Klammern ‚siehe
Anhang` schreibst.

Hinweise:

\begin{itemize}
\item
  Nummeriere alle Anhänge.
\item
  Gib allen Anhängen einen eindeutigen Titel.
\item
  Trenne die Anhänge, sodass jeder auf einer neuen Seite beginnt.
\end{itemize}

Mehr Informationen zum Anhang findest du unter:

\href{https://www.scribbr.de/apa-standard/anhang/}{Anhang laut
APA-Richtlinien}

Eidesstattliche Erklärung

Der Wortlaut der eidesstattlichen Erklärung wird dir meist von deinem
Fachbereich vorgegeben. Wenn das nicht der Fall ist, findest du hier ein
Beispiel:
\href{https://www.scribbr.de/masterarbeit/eidesstattliche-erklaerung-masterarbeit/}{Eidesstattliche
Erklärung Masterarbeit - Bedeutung + 7 Vorlagen}

Hiermit erkläre ich, dass ich die vorliegende Arbeit selbstständig
verfasst habe, dass ich sie zuvor an keiner anderen Hochschule und in
keinem anderen Studiengang als Prüfungsleistung eingereicht habe und
dass ich keine anderen als die angegebenen Quellen und Hilfsmittel
benutzt habe. Alle Stellen der Arbeit, die wörtlich oder sinngemäß aus
Veröffentlichungen oder aus anderweitigen fremden Äußerungen entnommen
wurden, sind als solche kenntlich gemacht.

Ort, Datum Unterschrift
